\documentclass{article}
\usepackage[utf8]{inputenc}
\begin{document}
\section{Exercise 1: Transitions!}

This exercise is intended to teach strategies for identifying problems with flow, transition and organization across writing.

\paragraph{} Your tasks for this exercise are the following:
\begin{itemize}
\item Read the following paragraph.
\item Edit as needed. How can you connect the ideas to improve the flow of the paragraph?
\item Feel free to combine sentences, split sentences or even remove or re-write them entirely. Talk to your neighbor! There are multiple approaches that you can take.
\item Re-read your final paragraph. Do the sentences connect and transition well?
\end{itemize}

Note: there is no need for additions or intense editing (grammar/spelling) - at least for this exercise! If you are running into any problems with Overleaf or have questions, please feel free to raise your hand and ask for help.

\subsection*{Sample Paragraph}

Extraction can simply be described as the “pulling out” of a substance from one phase to another. The technique involves placing the mixture (that is to undergo separation) in a separatory funnel. Two immiscible solvents are also added. The different components of the mixture will have differing solubilities. So the components will dissolve in either of the two solvents and form separate layers i.e. the mixture undergoes partitioning. Layers can be extracted and retained. Reactions (with either an acid or base) and procedures (such as filtration) may be necessary to obtain the purified product. 

An effective extraction has several potential factors. Certain factors must be considered e.g. the type and the quantity of solvent used, and the number of extractions carried out. These factors are necessary. They affect the physical process of partitioning. 


ewpage
\subsection*{Suggested Solution}

Extraction can simply be described as the “pulling out” of a substance from one phase to another.2 The technique involves placing the mixture, \textbf{which will undergo separation}, in a separatory funnel, \textbf{along with two immiscible solvents}. \textbf{Since} the components of the mixture \textbf{have different solubilities}, the components will dissolve in either of the two solvents and form separate layers i.e. the mixture undergoes partitioning.2 \textbf{These} layers can then be extracted, and retained. \textbf{Further} reactions (with either an acid or base) and procedures (such as filtration) may be necessary to obtain the purified product.

\textbf{To carry out an} effective extraction, certain factors must be considered e.g. the type and the quantity of solvent used, and the number of extractions carried out.1 These factors are necessary \textbf{as they} affect the physical process of partitioning. 

\end{document}
